\documentclass[10pt,a4paper]{amsart}
\usepackage[margin=19mm]{geometry}
\usepackage{amsmath,amssymb,mathtools}
\usepackage[scr=rsfs,cal=euler]{mathalfa}
\usepackage{xfrac}
\usepackage[unicode,hidelinks,bookmarks=false]{hyperref}
\newcommand{\ie}{i.e.\ }
\newcommand{\cf}{cf.\ }
\newcommand{\resp}{respectively\ }
\newcommand{\Subsection}{\subsection}
\providecommand{\nobreakdash}{\nobreak\hbox{-}}
\setlength{\emergencystretch}{2em}
\multlinegap0pt
\usepackage{amssymb}
\usepackage{mathtools}
\usepackage{bm}
\usepackage{enumerate}
\newtheorem{theorem}{Theorem}
\newtheorem{definition}{Definition}
\def\Identity{\mathrm{Id}} %
\def\Op#1{{\mathrm{#1}}}  % operator
\def\Spc#1{{\mathcal{#1}}}  % spaces
\def\V#1{{\boldsymbol{#1}}}         % vectors
\newcommand{\embedC}{\xhookrightarrow{}}
\title{Generalized Splines and Gaussian Processes}
\author{Michael Unser}
\date{Source: arXiv:2608.28446v1 (CC BY 4.0)}
\begin{document}
\maketitle

\noindent\emph{Source and licence.} Generalized Splines and Gaussian Processes. Version arXiv:2608.28446v1 (CC BY 4.0), distributed by arXiv under the Creative Commons Attribution 4.0 International licence, http://creativecommons.org/licenses/by/4.0/, as recorded in the arXiv licence metadata for this exact version and shown on the abstract page https://arxiv.org/abs/2608.28446v1. Full licence notice: \texttt{licenses/arxiv-2608.28446-CC-BY-4.0.txt}.\par\smallskip

\noindent\textit{Editorial scope.} Counted unit: the article's theorem Theo:conditionPos (source0.tex lines 467--483). The article's complete original proof of it is delivered with it (source0.tex lines 484--505, one \texttt{proof} environment of 22 lines). No mathematics is written, repaired or re-derived: the statement and the proof are the article's own lines, in the article's own notation, and no retained line is substituted or re-flowed.  Retained uncounted as the source-local context the proof needs: lines 350--366; lines 401--428.  Nothing was omitted from any retained span, so the package carries no omission record.  No retained line contains a citation, so the wrapper carries no bibliography and no bibliography entry.  No retained line contains a figure environment, an image-inclusion command or drawing code, so no image is delivered and the package declares no asset dependency.  Editorial formatting only, in addition: an AMS \texttt{amsart} preamble that restates the article's own declarations and macros used by the retained lines, namely the environments \texttt{theorem}, \texttt{definition} on separate counters, together with the standard packages those lines need.  The retained environments print in this document as Theorem 1, Definition 1 in order of appearance, whereas the article prints the same items with the numbering of its own full document.  The complete native source of the article as distributed in the arXiv e-print for this exact version is bundled unchanged as \texttt{sources/208789/source0.tex}.
\begin{theorem}[Finite-dimensional Hilbert subspaces of $\Spc S'$]
\label{Theo:FiniteDimHilbert}
Let $\V p=(p_n)$ be a basis of $\Spc N_{\V p}={\rm span}\{p_n\}_{n=1}^{N_0}\subset \Spc S'$. Then, there exits
 a vector of functionals $\V \phi=(\phi_n)$ with $\phi_1,\dots,\phi_{N_0}\in \Spc S$ such that
 $\langle p_n, \phi_m\rangle = \delta_{m-n}$ (biorthogonality).
 
The resulting biorthogonal system $(\V p, \V \phi)$ is called {\em universal} and it allows us to identify the dual pair of 
Hilbert spaces  $\Spc H_0=(\Spc N_{\V p},|\cdot|_{\V \phi}) \subset \Spc S'$ and
$\Spc H'_0=(\Spc N_{\V \phi},|\cdot|_{\V p}) \subset \Spc S$ equipped with the (semi-)inner
products
\begin{align}
\langle f_1,f_2\rangle_{\V \phi} = \sum_{n=1}^{N_0} \langle \phi_n,f_1 \rangle\langle \phi_n,f_2 \rangle\\ 
\langle \varphi_1,\varphi_2\rangle_{\V p} = \sum_{n=1}^{N_0} \langle p_n,\varphi_1 \rangle \langle p_n,\varphi_2 \rangle,
\end{align}
which are well-defined for any $f_1,f_2 \in \Spc S'$ and $\varphi_1,\varphi_2 \in \Spc S$.
The argument generalizes to $\Spc N_{\V p}\subset \Spc H$, and $\Spc N_{\V \phi}\subset \Spc H'$ where $\Spc H$ is any Hilbert space such that $\Spc S \embedC \Spc H \embedC \Spc S'$, which  enables the use of {\em non-universal} biorthogonal systems.
\end{theorem}

\begin{definition}[Positive-definite operator]
\label{Def:CondPositiveDef2}
Let  $\Op A$ be a continuous operator $\Spc S \to \Spc S'$ and $\Spc N_{\V p}$ a finite-dimensional subspace of $\Spc S'$ that is spanned by
$\V p=(p_1,\dots,p_{N_0})$.
The operator $\Op A$ is said to be
\begin{itemize}
\item symmetric or self-adjoint if, for all $\varphi_1,\varphi_2 \in \Spc S$,
$
\langle \Op A \varphi_1,\varphi_2\rangle=\langle \Op A \varphi_2,\varphi_1\rangle;
$
\item positive-definite if, for any  $\varphi \in \Spc S\backslash\{0\}
$,
$
\langle \Op A \varphi,\varphi\rangle > 0$;

\item positive-semi-definite if, for any  $\varphi \in \Spc S$,
$
\langle \Op A \varphi,\varphi\rangle \ge 0;
$\item $\V p$-conditionally positive-definite if
\begin{align}
\label{Eq:PositiveDefinite}
\langle \Op A \varphi,\varphi\rangle > 0
\mbox{ for all }\varphi \in  \Spc S_{\V p} \backslash\{0\},
\end{align}
where $\Spc S_{\V p}=\Spc S \cap \Spc N_{\V p}^\perp=\{\varphi \in \Spc S: \V p(\varphi)=\V 0\}$. 
%If \eqref{Eq:PositiveDefinite} holds for $\V p=\{0\}$, then $\Op A$ is said to be positive-definite or positive (for short).
\end{itemize}
\end{definition}

\begin{theorem}[Operator-based construction of Hilbert spaces]
\label{Theo:conditionPos}
Let $\Op A: \Spc S \to \Spc S'$ be a $\V p$-conditionally positive-definite operator.
% as in Definition \ref{Def:CondPositiveDef2}. % thereafter referred to as the Gram operator.
Then, for any given biorthogonal system $(\V p, \V \phi)$, the bilinear form
\begin{align}
\label{Eq:BPinnerA}
\langle f,g\rangle_{\Spc H'}=\langle(\Op A_{\V \phi} + \Op R_{\V p}) f,g\rangle
\end{align}
is a valid inner product on $\Spc S$. By completion, the latter specifies the Hilbert space $\Spc H'=\overline{(\Spc S,\|\cdot\|_{\Spc H'})}=\Spc H_{\Op A} \oplus \Spc N_{\V \phi}$, with
$\Spc H_{\Op A}=\{f \in \Spc H': \V p(f)=\V 0\}=\overline{(\Spc S_{\V p},\|\cdot\|_{\Spc H'})}$ the Hilbert space that is the orthogonal complement of $ \Spc N_{\V \phi}$ in $\Spc H'$. 

The predual of $\Spc H'$
is the Hilbert space $\Spc H=\Spc H'_{\Op A} \oplus \Spc N_{\V p}$ associated with the inner product $\langle f,g\rangle_{\Spc H}=\langle(\Op A^{-1} + \Op R_{\V \phi}) f,g\rangle$,
where $\Op A^{-1}$ is the unique operator such that $\Op A^{-1}|_{\Spc H'_{\Op A}}=\Op A_{\V \phi}^{-1}=\Op J_{\Spc H'_{\Op A}}: \Spc H'_{\Op A} \to \Spc H_{\Op A}$ (inverse of the Riesz map $\Op J_{\Spc H_{\Op A}}$) and $\Op A^{-1}|_{\Spc N_{\V p}}=0$ so that its null space is precisely $\Spc N_{\V p}$.

\end{theorem}

\begin{proof}[\small Proof]
{\small
We first select a universal system $\V \phi$ with $\phi_n \in \Spc S$, which ensures the continuity of
$\Op R_{\V \phi}: \Spc S' \to \Spc N_{\V \phi}$. Next, we define the seminorms $|\varphi|_{\Op A}=\sqrt{\langle \Op A_{\V \phi}\varphi,\varphi\rangle}$ and $|\varphi |_{\V p}=\sqrt{\langle \Op R_{\V p}\varphi,\varphi\rangle}$, with the latter reverting to the norm of the finite-dimensional Hilbert space $\Spc N_{\V \phi}$ for any $\varphi \in \Spc N_{\V \phi}$.
We then consider the direct-sum decomposition
$\Spc S=\Spc S_{\V p} \oplus \Spc N_{\V \phi}$ and observe that $|\varphi|_{\Op A}$ specifies a valid norm on
$\Spc S_{\V p}$ (from the definition of conditional positivity), while $|\phi|_{\Op A}$ vanishes for all $\phi \in \Spc N_{\V \phi}$ due to the reproduction condition $\Op R_{\V \phi}\Op R_\V p=\Identity$ on $\Spc N_{\V \phi}$. As for $|\varphi|_{\V p}$, we have exactly the reverse behavior with $|\varphi|_{\V p}=0$ for all $\varphi \in \Spc S_{\V p}$. These properties imply that
the augmented operator $\Op A_{\V \phi} + \Op R_{\V p}$ is positive-definite, with the induced norm
$\|\varphi\|_{\Spc H'}=\sqrt{|\varphi|^2_{\Op A}+|\varphi|^2_{\V p}}$ being compatible with the direct-sum decomposition.
This then allows us to specify the corresponding Hilbert spaces $\Spc H'=\overline{(\Spc S,\|\cdot\|_{\Spc H'})}$ and $\Spc H_{\Op A}=\overline{(\Spc S_{\V p},\|\cdot\|_{\Spc H'})}=\overline{(\Spc S_{\V p},|\cdot|_{\Op A})}$ by completion, with the direct-sum structure $\Spc H'=\Spc H_{\Op A} \oplus \Spc N_{\phi}$ being preserved in the process. 

An important point is that the (semi-)norm $|f|_{\Op A}$ on $\Spc H_{\Op A}=\{f \in \Spc H': \V p(f)=\V 0\}$ is not dependent on the choice of biorthogonal system, in the sense that
\begin{align}
\forall f \in \Spc H_{\Op A}: |f|^2_{\Op A}=\langle \Op A_{\V \phi} f, f\rangle=  \langle \Op A f, f\rangle \mbox{ for any } \V \phi \subset \Spc H' \mbox{ s.t. } \Op R_{\V \phi}\Op R_{\V p}=\Identity.
\end{align}
This allows us to extend the result to general biorthogonal systems subject to the constraint $\phi_1, \dots, \phi_{N_0} \in \Spc H'$, the bottom line being that these systems are all topologically equivalent,
with the underlying space $\Spc H'$ (as a set) remaining the same.

The second part of the theorem is obtained by duality, based on the property that 
$\Spc H=(\Spc H')'=(\Spc H_\Op A \oplus \Spc N_{\V \phi})'=\Spc H_\Op A' \oplus \Spc N'_{\V \phi}$. There, $\Spc N'_{\V \phi}=\Spc N_{\V p}$ (see Theorem \ref{Theo:FiniteDimHilbert}) and $\Spc H_\Op A'$ is a bona fide Hilbert space that is orthogonal to $\Spc N'_{\V \phi}$ (and, hence, to $\V p$) and whose Riesz map  $\Op J_{\Spc H'_\Op A}$ is the inverse of $\Op J_{\Spc H_\Op A}=\Op A: \Spc H_\Op A \to \Spc H'_\Op A$.
} % end small
\end{proof}

\end{document}
